%======================================================================
% SANDRA ECOSYSTEM - Diseño Revista Científica (Dos Columnas)
%======================================================================

\documentclass[11pt, spanish, letterpaper, twocolumn]{book}

\usepackage{fontspec}
\setmainfont{Inter}[Path = assets/fonts/, Extension = .otf, UprightFont = *-Regular]
\setsansfont{Inter}[Path = assets/fonts/, Extension = .otf, UprightFont = *-Medium]

\usepackage[
    top=2cm,
    bottom=2cm,
    left=2cm,
    right=2cm,
    columnsep=18pt
]{geometry}

\usepackage{graphicx}
\usepackage{wrapfig}
\usepackage{fancyhdr}
\usepackage{titlesec}
\usepackage{xcolor}
\usepackage{microtype}
\usepackage{tikz}
\usetikzlibrary{calc, shapes, positioning, arrows.meta}
\usepackage[most]{tcolorbox}
\usepackage{fontawesome5}
\usepackage{amsmath}
\usepackage{parskip}

%--------------------------------------------------------------------
% COLORES - Paleta elegante y científica
%--------------------------------------------------------------------
\definecolor{sandraDark}{HTML}{2C3E50}
\definecolor{sandraAccent}{HTML}{00897B}
\definecolor{sandraLight}{HTML}{E8F5E9}
\definecolor{sandraGray}{HTML}{ECEFF1}
\definecolor{sandraText}{HTML}{37474F}
\definecolor{isoBlue}{HTML}{1565C0}
\definecolor{isoGreen}{HTML}{2E7D32}
\definecolor{isoOrange}{HTML}{EF6C00}
\definecolor{isoPurple}{HTML}{7B1FA2}

%--------------------------------------------------------------------
% CONFIGURACIÓN DE PÁGINAS
%--------------------------------------------------------------------
\pagestyle{fancy}
\fancyhf{}
\fancyhead[RO]{\textcolor{sandraDark}{\itshape\rightmark}}
\fancyhead[LE]{\textcolor{sandraDark}{\itshape\leftmark}}
\fancyfoot[CE,CO]{\thepage}
\fancyfoot[RE]{\textcolor{sandraAccent}{\fontsize{8}{10}\selectfont\selectfont\itshape Sandra Ecosystem}}
\fancyfoot[LO]{\textcolor{sandraGray}{\fontsize{8}{10}\selectfont\selectfont\itshape Marco Teórico y Legal}}

\setlength{\headheight}{14pt}
\addtolength{\topmargin}{-2pt}

%--------------------------------------------------------------------
% TÍTULOS ELEgANTES
%--------------------------------------------------------------------
\titleformat{\chapter}[block]
{\normalfont\Huge\bfseries\color{sandraDark}\thispagestyle{empty}}{\thechapter}{0.5em}{}

\titleformat{\section}
{\normalfont\Large\bfseries\color{sandraDark}}{\thesection}{0.5em}{}

\titleformat{\subsection}
{\normalfont\large\bfseries\color{sandraAccent}}{\thesubsection}{0.5em}{}

\titlespacing*{\section}{0pt}{1.5ex plus 0.5ex minus 0.2ex}{0.5ex}
\titlespacing*{\subsection}{0pt}{1ex plus 0.3ex minus 0.1ex}{0.3ex}

%--------------------------------------------------------------------
% CAJAS DE INFORMACIÓN - Estilo científico limpio
%--------------------------------------------------------------------
\newtcolorbox{notacientifica}[1][]{
    colback=sandraGray,
    colframe=sandraAccent,
    boxrule=0.5pt,
    rounded corners,
    sharp corners,
    size=fbox,
    after skip=0.5ex,
    #1
}

\newtcolorbox{cajaestandar}[1][]{
    colback=sandraLight,
    colframe=isoBlue,
    boxrule=1pt,
    rounded corners,
    size=fbox,
    after skip=0.5ex,
    #1
}

\newtcolorbox{cajaadvertencia}[1][]{
    colback=#FFF3E0,
    colframe=isoOrange,
    boxrule=1pt,
    rounded corners,
    size=fbox,
    after skip=0.5ex,
    #1
}

%--------------------------------------------------------------------
% FIGURAS CON caption MEJORADO
%--------------------------------------------------------------------
\captionsetup[figure]{font=small,labelfont=bf,labelsep=period}
\captionsetup[table]{font=small,labelfont=bf,labelsep=period}

%--------------------------------------------------------------------
% DROP CAP ELEGANTE
%--------------------------------------------------------------------
\usepackage{lettrine}
\lettrinepreludesearch={l}
\lettrineafterlen=2pt
\lettrinelinenums=true
\lettrinefont{\fontsize{36}{40}\selectfont\bfseries\color{sandraDark}}

%--------------------------------------------------------------------
% ENTORNOS
%--------------------------------------------------------------------
\newenvironment{concepto}
{\begin{trivlist}\item[\hskip\labelsep{\bfseries\color{sandraAccent}\faIcon{lightbulb}}]\itshape}
{\end{trivlist}}

%--------------------------------------------------------------------
% COMANDOS AUXILIARES
%--------------------------------------------------------------------
\newcommand{\tecnologia}[2]{%
    \begin{tcolorbox}[colback=white, colframe=sandraGray, boxrule=0.3pt, 
        before skip=0.3ex, after skip=0.3ex, left=2pt, right=2pt, top=2pt, bottom=2pt]
        \textbf{\color{sandraDark}#1}\quad\color{sandraText}#2
    \end{tcolorbox}%
}

\newcommand{\norma}[2]{%
    \tikz[baseline=(char.base)]{%
        \node[circle, fill=isoBlue, text=white, font=\bfseries\tiny, inner sep=1pt] (char) {#1};
    }\,#2%
}

\usepackage[spanish]{babel}
\addto\captionsspanish{\def\tablename{Tabla}}

\begin{document}

%======================================================================
% PÁGINA DE TÍTULO DE CAPÍTULO - Estilo revista científica
%======================================================================
\thispagestyle{empty}
\null\vfill
\begin{center}
    % Número de capítulo estilo revista
    \begin{tikzpicture}[remember picture, overlay]
        \node[anchor=north east] at ([xshift=-1cm, yshift=-1cm] current page.north east) {
            {\fontsize{72}{80}\selectfont\bfseries\color{sandraAccent!30}02}
        };
    \end{tikzpicture}
    
    {\fontsize{11}{13}\selectfont\color{sandraAccent}\bfseries\uppercase{Capítulo 2}\par}
    \vspace{0.8cm}
    {\fontsize{28}{34}\selectfont\bfseries\color{sandraDark}Marco Teórico\par}
    \vspace{0.3cm}
    {\fontsize{12}{15}\selectfont\color{sandraText}\itshape Fundamentos tecnológicos y arquitectura de sistemas empresariales\par}
    \vspace{1cm}
    \rule{0.25\linewidth}{1pt}
    \vspace{0.5cm}
    {\fontsize{10}{12}\selectfont\color{sandraGray}Carlos Peña}
\end{center}
\vfill
\clearpage

%--------------------------------------------------------------------
% ESTADO DEL ARTE
%--------------------------------------------------------------------
\section{Estado del Arte}

\lettrine{E}{l} middleware empresarial ha evolucionado significativamente desde sus orígenes hasta convertirse en plataformas complejas de integración. El mercado ofrece soluciones maduras como IBM Integration Bus, MuleSoft, Apache Camel y Spring Integration, cada una con fortalezas específicas para diferentes escenarios de uso.

Sin embargo, la mayoría de estas soluciones fueron diseñadas para mercados desarrollados, ignorando las necesidades específicas de regiones con infraestructura limitada y contextos regulatorios distintos. Esta brecha crea una oportunidad clara para soluciones localmente desarrolladas que aborden estos desafíos particulares.

\begin{notacientifica}
    \textbf{Nota del Autor}: El desarrollo de middleware propio representa una oportunidad estratégica para garantizar la soberanía tecnológica en entornos gubernamentales y empresariales críticos.
\end{notacientifica}

%--------------------------------------------------------------------
% FUNDAMENTOS DE MIDDLEWARE
%--------------------------------------------------------------------
\section{Fundamentos de Middleware}

El middleware es un software que actúa como puente entre aplicaciones y el sistema operativo, facilitando la comunicación e integración entre sistemas heterogéneos. Un Enterprise Service Bus (ESB) representa la evolución hacia una arquitectura basada en servicios, proporcionando un medio centralizado para la integración.

\begin{wrapfigure}{R}{0.45\linewidth}
    \centering
    \begin{tikzpicture}[
        node distance=0.6cm,
        rect/.style={rectangle, draw, rounded corners=2pt, minimum height=0.5cm, font=\fontsize{7}{8}\selectfont},
        arrow/.style={->, >=stealth, thin}
    ]
        \node[rect, fill=sandraLight] (a1) {Aplicación 1};
        \node[rect, fill=sandraLight] (a2) {Aplicación 2};
        \node[rect, fill=sandraLight] (a3) {Aplicación N};
        
        \node[rect, fill=sandraAccent!30, text centered, text width=2cm, minimum height=0.8cm] (esb) at ([yshift=-1.2cm] a2) {Enterprise\\Service Bus};
        
        \node[rect, fill=sandraGray, text centered] (so) at ([yshift=-1.2cm] esb) {Sistema Operativo};
        
        \draw[arrow] (a1) -- (esb);
        \draw[arrow] (a2) -- (esb);
        \draw[arrow] (a3) -- (esb);
        \draw[arrow] (esb) -- (so);
    \end{tikzpicture}
    \caption{Arquitectura básica de middleware}
    \label{fig:middleware}
\end{wrapfigure}

Sandra Server implementa estos conceptos utilizando Go, proporcionando un ESB de alto rendimiento capaz de manejar miles de conexiones simultáneas con consumo mínimo de recursos. Esta implementación permite:

\begin{itemize}
    \item Transformación de mensajes entre diferentes formatos
    \item Enrutamiento basado en reglas de negocio
    \item Composición de servicios complejos
    \item Manejo de transacciones distribuidas
\end{itemize}

%--------------------------------------------------------------------
% ARQUITECTURA DE MICROSERVICIOS
%--------------------------------------------------------------------
\section{Arquitectura de Microservicios}

La arquitectura de microservicios construye aplicaciones como servicios pequeños, autónomos y débilmente acoplados. Cada microservicio es responsable de una funcionalidad específica y puede desplegarse independientemente, permitiendo escalabilidad granular y tolerancia a fallos.

\begin{cajaestandar}
    \textbf{Ecosistema Sandra}: El ecosistema implementa esta arquitectura distribuida con cuatro componentes principales que se comunican mediante APIs REST y gRPC.
\end{cajaestandar}

\begin{table}[H]
    \centering
    \caption{Componentes del ecosistema Sandra}
    \begin{tabular}{lp{6cm}}
        \toprule
        \textbf{Componente} & \textbf{Descripción} \\
        \midrule
        \textbf{Server} & Middleware central de comunicaciones \\
        \textbf{Consola} & Interfaz administrativa web \\
        \textbf{SDC} & Aplicación de escritorio multiplataforma \\
        \textbf{Sentinel} & Procesamiento especializado de eventos \\
        \bottomrule
    \end{tabular}
    \label{tab:componentes}
\end{table}

%--------------------------------------------------------------------
% SEGURIDAD ZERO TRUST
%--------------------------------------------------------------------
\section{Modelo de Seguridad Zero Trust}

El modelo Zero Trust elimina la confianza implícita, verificando cada solicitud independientemente de su origen en la red. Este paradigma asume que ningún usuario o dispositivo es confiable por defecto, requiriendo autenticación y autorización continuas.

\norma{ISO}{27001}: La familia de normas ISO/IEC 27000 proporciona un marco sistemático para gestionar la seguridad de la información, implementando controles que abarcan aspectos organizacionales, personales, físicos y tecnológicos.

\begin{cajaadvertencia}
    \textbf{Consideración importante}: La implementación de Zero Trust requiere un cambio cultural organizacional, no solo implementación técnica.
\end{cajaadvertencia}

Los principios fundamentales incluyen: verificación explícita de usuarios y dispositivos mediante autenticación multifactor; acceso con mínimo privilegio, otorgando solo los permisos necesarios; y verificación continua de seguridad mediante monitoreo de sesiones y detección de anomalías.

Sandra Ecosystem implementa estos principios en todas las capas: cada solicitud de API es autenticada y autorizada, las comunicaciones usan TLS mutuo, y las sesiones tienen duración limitada.

%--------------------------------------------------------------------
% TECNOLOGÍAS
%--------------------------------------------------------------------
\section{Tecnologías Emergentes}

El ecosistema utiliza tecnologías de última generación, seleccionadas por su rendimiento, seguridad y comunidad activa:

\begin{figure}[H]
    \centering
    \begin{tikzpicture}[
        node distance=0.4cm,
        box/.style={rectangle, draw, rounded corners=3pt, minimum width=1.4cm, minimum height=0.6cm, font=\fontsize{6}{7}\selectfont, align=center},
        arrow/.style={->, >=stealth, thin}
    ]
        \node[box, fill=sandraLight] (go) {Go 1.24};
        \node[box, fill=sandraLight, right=0.3cm of go] (rust) {Rust};
        \node[box, fill=sandraLight, right=0.3cm of rust] (angular) {Angular};
        \node[box, fill=sandraLight, right=0.3cm of angular] (tauri) {Tauri};
        \node[box, fill=sandraLight, right=0.3cm of tauri] (grpc) {gRPC};
        
        \node[box, fill=sandraGray, below=0.3cm of go] (goD) {Rendimiento\\Concurrencia};
        \node[box, fill=sandraGray, below=0.3cm of rust] (rustD) {Seguridad\\Memoria};
        \node[box, fill=sandraGray, below=0.3cm of angular] (angularD) {Frontend\\Enterprise};
        \node[box, fill=sandraGray, below=0.3cm of tauri] (tauriD) {Escritorio\\Ligero};
        \node[box, fill=sandraGray, below=0.3cm of grpc] (grpcD) {API\\Rápida};
        
        \draw[arrow] (go) -- (goD);
        \draw[arrow] (rust) -- (rustD);
        \draw[arrow] (angular) -- (angularD);
        \draw[arrow] (tauri) -- (tauriD);
        \draw[arrow] (grpc) -- (grpcD);
    \end{tikzpicture}
    \caption{Tecnologías del ecosistema}
    \label{fig:tecnologias}
\end{figure}

\begin{notacientifica}
    La combinación de Go para el backend y Tauri para el frontend proporciona un equilibrio óptimo entre rendimiento y consumo de recursos, resultando en aplicaciones que pueden ejecutarse en hardware modesto.
\end{notacientifica}

\clearpage

%======================================================================
% CAPÍTULO 3: MARCO LEGAL - Página de título
%======================================================================
\thispagestyle{empty}
\null\vfill
\begin{center}
    \begin{tikzpicture}[remember picture, overlay]
        \node[anchor=north east] at ([xshift=-1cm, yshift=-1cm] current page.north east) {
            {\fontsize{72}{80}\selectfont\bfseries\color{sandraAccent!30}03}
        };
    \end{tikzpicture}
    
    {\fontsize{11}{13}\selectfont\color{sandraAccent}\bfseries\uppercase{Capítulo 3}\par}
    \vspace{0.8cm}
    {\fontsize{28}{34}\selectfont\bfseries\color{sandraDark}Marco Legal\par}
    \vspace{0.3cm}
    {\fontsize{12}{15}\selectfont\color{sandraText}\itshape Normativas y cumplimiento legal en Venezuela\par}
    \vspace{1cm}
    \rule{0.25\linewidth}{1pt}
    \vspace{0.5cm}
    {\fontsize{10}{12}\selectfont\color{sandraGray}Carlos Peña}
\end{center}
\vfill
\clearpage

%--------------------------------------------------------------------
% LEY DE INFOGOBIERNO
%--------------------------------------------------------------------
\section{Ley de Infogobierno}

La Ley de Infogobierno establece normas para el uso de tecnologías de información en el sector público venezolano, promoviendo la Soberanía Tecnológica y el uso de Software Libre. Esta legislación representa un avance significativo hacia la reducción de la dependencia tecnológica extranjera.

\begin{table}[H]
    \centering
    \caption{Cumplimiento con la Ley de Infogobierno}
    \begin{tabular}{lp{5.5cm}}
        \toprule
        \textbf{Requisito} & \textbf{Implementación en Sandra} \\
        \midrule
        Software Libre & Todas las tecnologías son código abierto \\
        Estándares Abiertos & REST, gRPC, JSON, Protocol Buffers \\
        Interoperabilidad & APIs documentadas, adaptadores modulares \\
        Seguridad & Cifrado AES-256, JWT, MFA, TLS 1.3 \\
        \bottomrule
    \end{tabular}
    \label{tab:infogobierno}
\end{table}

%--------------------------------------------------------------------
% LEY ESPECIAL CONTRA DELITOS INFORMÁTICOS
%--------------------------------------------------------------------
\section{Ley Especial contra los Delitos Informáticos}

Esta ley establece conductas penalizadas relacionadas con el uso indebido de sistemas informáticos. Sandra Ecosystem implementa controles que ayudan a prevenir y detectar estas conductas:

\begin{wrapfigure}{R}{0.4\linewidth}
    \centering
    \begin{tikzpicture}[
        node distance=0.3cm,
        box/.style={rectangle, draw, rounded corners=2pt, minimum height=0.4cm, font=\fontsize{6}{7}\selectfont, fill=sandraLight},
        arrow/.style={->, >=stealth, thin}
    ]
        \node[box] (usr) {Usuario};
        \node[box, below=0.3cm of usr] (auth) {Autenticación};
        \node[box, below=0.3cm of auth] (perm) {Permisos};
        \node[box, below=0.3cm of perm] (rec) {Recurso};
        
        \draw[arrow] (usr) -- (auth);
        \draw[arrow] (auth) -- (perm);
        \draw[arrow] (perm) -- (rec);
    \end{tikzpicture}
    \caption{Control de acceso}
    \label{fig:acceso}
\end{wrapfigure}

\begin{itemize}
    \item \textbf{Acceso no autorizado}: Autenticación multifactor (MFA)
    \item \textbf{Interceptación de datos}: Cifrado en tránsito y en reposo
    \item \textbf{Falsificación de datos}: Firmas digitales y logs de auditoría
    \item \textbf{Destrucción de datos}: Sistemas de backup y redundancia
\end{itemize}

%--------------------------------------------------------------------
% ESTÁNDARES ISO
%--------------------------------------------------------------------
\section{Estándares Internacionales}

\norma{ISO}{27001}: Sistema de Gestión de Seguridad de la Información (SGSI). Proporciona un enfoque sistemático para gestionar información sensible.

\norma{ISO}{27002}: Controles de seguridad de la información. Incluye 93 medidas organizadas en cuatro categorías: organizacionales, personales, físicos y tecnológicos.

\norma{NIST}{SP 800-53}: Controles de seguridad y privacidad del Instituto Nacional de Estándares y Tecnología de Estados Unidos.

\begin{cajaestandar}
    Sandra Ecosystem implementa controles alineados con ISO 27001, NIST y mejores prácticas de la industria para garantizar la confidencialidad, integridad y disponibilidad de la información.
\end{cajaestandar}

%--------------------------------------------------------------------
% LICENCIAMIENTO ESTRATÉGICO
%--------------------------------------------------------------------
\section{Licenciamiento Soberano}

La elección de licencias en el Proyecto Sandra no es solo legal, es estratégica. Se priorizan licencias que garantizan la autonomía tecnológica y la protección de la propiedad intelectual en entornos críticos.

\begin{cajaadvertencia}
    \textbf{Nota estratégica}: Las licencias permisivas como Apache 2.0 y MIT permiten que innovaciones basadas en código abierto puedan ser privatizadas y clasificadas, un requerimiento para el sector defensa.
\end{cajaadvertencia}

\begin{table}[H]
    \centering
    \caption{Análisis de licenciamiento}
    \begin{tabular}{lp{5.5cm}}
        \toprule
        \textbf{Licencia} & \textbf{Implicación} \\
        \midrule
        \textbf{Apache 2.0 / MIT} & Recomendado: Permite clasificar módulos críticos \\
        \textbf{GPL v3 / AGPL} & Restringido: Copyleft fuerte prohíbe cierre \\
        \bottomrule
    \end{tabular}
    \label{tab:licencias}
\end{table}

\begin{notacientifica}
    Si desarrollas un módulo de guiado basado en un algoritmo abierto, una licencia como Apache 2.0 permite que el software final sea cerrado y clasificado sin comprometer el cumplimiento legal.
\end{notacientifica}

\end{document}
